\section{Discusi\'on}
\label{sec:discusion}

\subsection{Lo que el \'indice s\'i mide}

El patr\'on geogr\'afico del \'indice (Figura \ref{fig:mapa}) es consistente con la sismolog\'ia y la hidrolog\'ia conocidas de M\'exico: las presas de mayor peligro s\'ismico se concentran en la zona de subducci\'on de la placa de Cocos a lo largo de la costa del Pac\'ifico (Chiapas, Oaxaca, Guerrero, Michoac\'an, Colima), mientras que el altiplano central --geol\'ogicamente m\'as estable-- domina el extremo de baja prioridad. La aparici\'on de presas de la frontera de Baja California y de la zona metropolitana de Puebla en el top 10 --pese a su sismicidad moderada-- ilustra precisamente el valor de combinar peligro con exposici\'on: son presas donde el evento detonante es menos probable, pero donde una falla afectar\'ia a una poblaci\'on much\'isimo mayor. Ninguno de los dos factores por separado habr\'ia identificado ese patr\'on.

\subsection{Lo que el \'indice todav\'ia no mide}

Este trabajo mide \emph{peligro} y \emph{exposici\'on}, pero no \emph{vulnerabilidad} --la tercera dimensi\'on cl\'asica del riesgo (Secci\'on \ref{sec:marco_teorico})--. Una presa puede tener baja prioridad en este \'indice por estar en una zona de peligro moderado y baja densidad poblacional cercana, y a\'un as\'i representar un riesgo alto si su estructura est\'a severamente deteriorada. Ese dato --altura de cortina, a\~no de construcci\'on, capacidad del vertedor, estado f\'isico y fecha de \'ultima inspecci\'on-- es precisamente el que CONAGUA no publica para la mayor parte del inventario (Secci\'on \ref{sec:introduccion}), y el que se solicit\'o por separado v\'ia transparencia gubernamental. Hasta que esa informaci\'on est\'e disponible, este \'indice debe leerse como una priorizaci\'on de \emph{d\'onde mirar primero}, no como un diagn\'ostico de seguridad estructural.

Una segunda limitaci\'on es de escala: el an\'alisis cubre las 210 presas que CONAGUA ya monitorea activamente, que son estructuralmente las mejor vigiladas del pa\'is --no las \(\sim\)1{,}000 sin inspecci\'on reciente que motivan este trabajo (Secci\'on \ref{sec:introduccion})--, ni los \(\sim\)8{,}000 bordos no registrados. El valor de este resultado es metodol\'ogico: demuestra que el \'indice es calculable y produce un ranking geogr\'aficamente coherente sobre un subconjunto verificable, listo para aplicarse al inventario completo en cuanto est\'e disponible.

Una tercera limitaci\'on es de precisi\'on de cada factor individual. El PGA se calcula para condiciones de roca de referencia y no incorpora amplificaci\'on por sitio (Secci\'on \ref{sec:metodologia_sismico}); la poblaci\'on cercana es un radio en l\'inea recta, no una zona de inundaci\'on real trazada por un modelo de elevaci\'on digital (Secci\'on \ref{sec:metodologia_poblacion}); y la lluvia de dise\~no asume que la distribuci\'on Gumbel ajustada sobre 33 a\~nos de registro es representativa del r\'egimen extremo de cada sitio, sin correcci\'on por tendencias de cambio clim\'atico. Cada una de estas simplificaciones es razonable para un \emph{screening} de primera pasada, pero ninguna sustituye un estudio de sitio o un modelo hidrol\'ogico de detalle para una presa espec\'ifica que el \'indice se\~nale como prioritaria.

\subsection{Sobre la dependencia de una solicitud de transparencia}

Que el paso metodol\'ogico m\'as importante pendiente de este trabajo --incorporar vulnerabilidad estructural real-- dependa de una solicitud de acceso a informaci\'on p\'ublica, y no de una limitaci\'on t\'ecnica, es en s\'i mismo un hallazgo relevante: refuerza emp\'iricamente el reclamo de CONAGUA de que el vac\'io de supervisi\'on de presas peque\~nas es, ante todo, un problema de recursos y de transparencia, no de falta de metodolog\'ias disponibles para atenderlo.
